% Part: set-theory
% Chapter: ordinals
% Section: replacement

\documentclass[../../../include/open-logic-section]{subfiles}

\begin{document}

\olfileid{sth}{ordinals}{replacement}
\olsection{প্রতিস্থাপন}

\olref[sth][ordinals][vn]{sec}-এ ক্রমসংখ্যা প্রবর্তনের পক্ষে
আমরা বলেছিলাম, এগুলিকে ক্রমপ্রকারের, অর্থাৎ
সুক্রমবিন্যাসের আদর্শ প্রতিনিধিরূপে ব্যবহার করা যেতে পারে।
তা করতে হলে প্রমাণ করতে হবে যে
\emph{প্রত্যেক সুক্রমবিন্যাস কোনো একটি ক্রমসংখ্যার সঙ্গে সমরূপ}।
তখন $\tuple{A, <} \isomorphic \alpha$ এমন ক্রমসংখ্যা
$\alpha$ দিয়ে $\ordtype{A, <}$ সংজ্ঞায়িত করা যাবে।

% BN-SRC-437: source omits a preposition and repeats a verb in the claim about previously introduced axioms.
কিন্তু এ পর্যন্ত প্রবর্তিত স্বতঃসিদ্ধগুলি \emph{শুধু} ব্যবহার করে
আমরা চাওয়া ফলটি প্রমাণ করতে \emph{পারি না}। কেন পারি না,
তা \olref[replacement][strength]{sec}-এ দেখব; আপাতত এই
সীমাবদ্ধতাটিই গুরুত্বপূর্ণ। আমাদের একটি নতুন ধারণা দরকার:

\begin{axiom}[প্রতিস্থাপন ছক]
যেকোনো সূত্র $\phi(x, y)$-এর জন্য নিচেরটি একটি স্বতঃসিদ্ধ:
\begin{quote}
	যেকোনো $A$-এর জন্য, যদি $(\forall x \in A)\lexists![y][\phi(x,y)]$ হয়,
	তবে $\Setabs{y}{(\exists x \in A)\phi(x,y)}$ বিদ্যমান।
\end{quote}
\end{axiom}
\noindent
পৃথকীকরণের মতো এটিও একটি ছক: অসংখ্য ভিন্ন $\phi$-এর
প্রতিটির জন্য এটি একটি করে স্বতঃসিদ্ধ দেয়। একই বক্তব্য
নিচের মতোও লেখা যায়, আর সচরাচর এভাবেই লেখা হয়:

\begin{defish}
``$B$'' নেই এমন যেকোনো সূত্র $\phi(x,y)$-এর জন্য নিচেরটি একটি স্বতঃসিদ্ধ:
\[
\forall A[(\forall x \in A)\lexists![y][\phi(x,y)] \lif \exists B\forall y (y \in B \liff (\exists x \in A)\phi(x,y))]
\]
\end{defish}

প্রথম দেখায় অবশ্য সূত্রটি বেশ জটিল লাগে। প্রতিস্থাপন থেকে
পাওয়া নিচের সহজ পরিণামটি আমাদের মূল ধারণা আরও
\emph{স্পষ্টভাবে} প্রকাশ করে:\footnote{পদ হলো এমন একটি রাশি, যা
তার মুক্ত চলকগুলির যেকোনো মান নির্ধারণে ঠিক একটি বস্তুকে নির্দেশ
করে। যেমন, ``$\emptyset$'' পদটি শূন্য সেটকে নির্দেশ করে;
``$\{x\}$'' পদটি $x$-এর একক-সদস্যবিশিষ্ট সেটকে নির্দেশ করে
($x$ যাই হোক); আর ``$x \cup y$'' পদটি $x$ ও $y$-এর
সংযোগকে নির্দেশ করে (তারা যাই হোক)।}

\begin{cor}
যেকোনো পদ $\tau(x)$ এবং যেকোনো সেট $A$-এর জন্য
নিচের সেটটি বিদ্যমান:
\[
	\Setabs{\tau(x)}{x \in A} = \Setabs{y}{(\exists x \in A)y = \tau(x)}.
\]
\end{cor}

\begin{proof}
$\tau$ একটি \emph{পদ} বলে $\forall x \lexists![y][\tau(x) = y]$।
বিশেষত, $(\forall x \in A)\lexists![y][\tau(x) = y]$। তাই
প্রতিস্থাপন অনুযায়ী $\Setabs{y}{(\exists x \in A)\tau(x) = y}$
বিদ্যমান।
\end{proof}
\noindent
এ থেকেই ``প্রতিস্থাপন'' নামটির তাৎপর্য বোঝা যায়: একটি
সেট $A$ দেওয়া থাকলে $A$-এর প্রত্যেক সদস্যের জায়গায় $\tau$-র অধীনে তার
প্রতিচ্ছবি বসিয়ে নতুন সেট $\Setabs{\tau(x)}{x \in A}$
গঠন করা যায়। ফাংশনের অধীনে কোনো সেটের প্রতিচ্ছবির
সংকেতলিপি অনুসরণ করে $\Setabs{\tau(x)}{x \in A}$-এর
জন্য $\funimage{\tau}{A}$-ও লিখতে পারি।

তবে এখানে একটি জরুরি পার্থক্য আছে: $\tau$ একটি \emph{পদ}।
\olref[sfr][fun][rel]{sec}-এর অর্থে, অর্থাৎ ক্রমজোড়ের
বিশেষ একটি সেট হিসেবে, এটি \emph{ফাংশন} হওয়া জরুরি নয়;
ফাংশনের নাম হওয়াও জরুরি নয়। কারণ $f$ ওই অর্থে একটি
ফাংশন হলে $\funimage{f}{A} = \Setabs{f(x)}{x \in
A}$ হলো $\ran{f}$-এর একটি বিশেষ উপসেট, যার অস্তিত্ব
শুধু $\Zminus$-এর স্বতঃসিদ্ধগুলি দিয়েই নিশ্চিত করা যায়।\footnote{যেমন
$\Setabs{y \in \bigcup \bigcup f}{(\exists x \in A)y =
f(x)}$ বিবেচনা করো।} কিন্তু প্রতিস্থাপন আমাদের
স্বতঃসিদ্ধসমূহে একটি \emph{শক্তিশালী} সংযোজন; তা
\olref[sth][replacement][]{chap}-এ দেখব।

\end{document}
